\chapter*{Abstract}
\addcontentsline{toc}{chapter}{Abstract}

This specification defines the Human-in-the-Loop (HITL) Requirements
Traceability framework for Standard Chartered Bank's Trial Balance (TB)
review process. It provides the complete audit trail from source documents
(business cases, detailed operating instructions, BPMN process maps, and
workbooks) through machine-readable extraction, structured business
requirements, formal specifications, and validating JSON schemas.

The document serves as the \textbf{human validation checkpoint} for the
automated requirements extraction and traceability process, enabling
reviewers to verify the accuracy of automated extractions and establish
a clear audit trail from source to implementation.

Four orthogonal workstreams are traced in full:

\begin{itemize}
  \item a \textbf{source document inventory} cataloguing 6 source files
        totalling 137~MB, including business cases, DOI, BPMN process map,
        and HKG TB/PL review workbooks;
  \item an \textbf{extraction layer} of 7 machine-readable outputs
        (markdown, JSON, YAML, CSV) validated against YAML frontmatter
        schemas;
  \item a \textbf{requirements layer} of 27 business requirements
        categorised across executive summary, scope, process steps,
        AI augmentation, and BSS risk assessment;
  \item a \textbf{specification and schema layer} comprising 21
        specification elements and 6 JSON schemas with field-level
        traceability back to requirements.
\end{itemize}

\paragraph{Reading order.} \Cref{ch:overview} introduces the HITL framework
and traceability model; \cref{ch:source-documents} catalogues source
documents; \cref{ch:extraction-layer} covers machine-readable extraction;
\cref{ch:business-requirements} details requirement derivation;
\cref{ch:specification-mapping} maps specifications to requirements;
\cref{ch:schema-mapping} provides schema field traceability;
\cref{ch:bidirectional-traceability} presents the full forward and
backward matrices; \cref{ch:gap-analysis} identifies traceability gaps;
\cref{ch:human-review-signoff} captures review checkpoints and approvals.

\paragraph{Conventions.} Source document references appear as
\srcref{TB-BC-001}. Business requirement IDs appear as \reqref{TB-REQ-ES01}.
Specification elements appear as \specref{TB-STEP-003}. Schema references
appear as \schemaref{variance-record}. Gap references appear as
\gapref{G-001}. File paths are typeset in \texttt{monospace} and are
relative to the repository root.

\paragraph{Companion files.} The machine-readable companion to this
specification is \texttt{docs/tb/machine-readable/hitl-traceability-manifest.yaml},
which enables automated validation of all traceability links documented
in this specification.

\tableofcontents